\chapter{Симметрийный допуск изотропии шумового модуля}
\label{ch:v8-noise-isotropy-symmetry-admission-gate}

\section{Проверка сильного чтения общего следа}

Предыдущий гейт построил базис-независимый оператор Казимира 19-мерного шумового
факторпространства. Для физической единственности его изотропия должна следовать из
действующей симметрии, а сам шумовой модуль должен быть инвариантен.
Проверим это непосредственно для представленной группы
$SU(3)\times SU(2)\times U(1)$ и оператора цепной степени.

\section{Первая поправка: выбранный модуль не замкнут полной группой}

Семь выбранных направлений переноса состоят из связывающей инцидентности и
двух цветных межсекторных семейств. При двенадцати случайных конечных калибровочных
преобразованиях максимальная норма выхода из этого подпространства равна
\begin{equation}
 2.21094673.
 \label{eq:v8-selected-noise-gauge-leakage}
\end{equation}
Следовательно, прежний 19-мерный оператор Казимира не является оператором Казимира полной калибровочной
группы.

Причина физически прозрачна: фон $A_0$ сам имеет нетривиальную калибровочную орбиту.
Ранг касательного отображения орбиты равен трём, его стабилизатор в
двенадцатимерной представленной калибровочной алгебре имеет размерность девять.
Под этим стабилизатором выбранный 19-мерный модуль действительно замкнут с
остатками порядка $10^{-15}$. Поэтому прежний результат сохраняется как
каноническая конструкция в фиксированном вакууме, но не как полная
калибровочно ковариантная конструкция до выбора вакуума.

\section{Минимальное полное калибровочное замыкание}

Комплексный линейная оболочка двадцати тяжёлых вещественных/мнимых вариаций имеет размерность
десять и уже инвариантен. Калибровочная орбита связывающей инцидентности имеет
размерность пять. Их сумма прямая и даёт
\begin{equation}
 \dim_{\mathbb C}\mathcal N_{\rm tr}^{\rm gauge}=5+10=15.
 \label{eq:v8-full-gauge-transfer-dimension}
\end{equation}
В сравнении с выбранным межсекторным модулем необходимо вернуть четыре комплексных
слабых направления семейств $LLXR,YLeR$ и четыре недостающих направления
орбиты $A_0$. Вместе с двенадцатью калибровочными шумами минимальное полное факторпространство
имеет размерность
\begin{equation}
 \boxed{\dim_{\mathbb C}\mathcal N_{\rm full}=15+12=27.}
 \label{eq:v8-full-gauge-noise-dimension}
\end{equation}
На этом расширении все двенадцать конечных калибровочных проверок замыкаются;
максимальные остатки инвариантности и унитарности меньше $6.4\cdot10^{-15}$.

\section{Коммутант не скалярен в обеих постановках}

В ветви фиксированного вакуума коммутант стабилизатора и цепной степени на
19-мерном модуле имеет комплексную размерность десять. В полной калибровочной ветви
на 27-мерном модуле его размерность равна шестнадцати:
\begin{equation}
 \dim\operatorname{Comm}_{19}=10,
 \qquad
 \dim\operatorname{Comm}_{27}=16.
 \label{eq:v8-noise-symmetry-commutants}
\end{equation}
Значит, ни одна физически честная ветвь не делает шумовое представление
неприводимым.

Уже два проекторных блока степени дают допустимое семейство метрик
\begin{equation}
 G_\eta=\eta P_{\rm tr}+P_{\rm gauge},
 \qquad \eta>0.
 \label{eq:v8-symmetry-invariant-noise-family}
\end{equation}
Все тринадцать значений $10^{-3}\leq\eta\leq10^3$ сохранили KMS-симметрию,
примитивность и одномерную неподвижную алгебру; щель менялась от
$3.8970\cdot10^{-4}$ до $389.7003$. Симметрия не выбирает $\eta=1$.

Формальная группа $U(27)$ сделала бы коммутант скалярным, но типичное её
преобразование не сохраняет цепную степень: норма коммутатора равна
$7.1541$. Такое расширение смешивает нулевые и боровские частоты два и не
может быть объявлено физической симметрией.

\begin{theorem}[Две симметрийные ветви шумового модуля]
После фиксации связывающего фона 19-мерное шумовое факторпространство ковариантно только
относительно девятимерного стабилизатора $A_0$. До фиксации фона минимальное
полное калибровочное замыкание имеет размерность 27. Коммутанты обеих ветвей
нетривиальны, поэтому физическая симметрия не вынуждает следовую изотропию.
\end{theorem}

\begin{nogocriterion}[Запрет полного калибровочного статуса 19-мерного Казимир]
Нельзя использовать веса прежнего 19-мерного оператора Казимира как полностью калибровочно-
ковариантные до выбора вакуума. Они относятся к стабилизаторной ветви.
Полная ветвь требует 27-мерного шумового факторпространства и повторной проверки его
родительского гессиана. Нельзя также постулировать $U(27)$: она нарушает
точную цепную степень.
\end{nogocriterion}

\begin{protocol}
\begin{enumerate}
 \item выбранный модуль переноса: \textbf{размерность 7};
 \item его полная калибровочная орбита: \textbf{размерность 11};
 \item стабилизатор $A_0$: \textbf{размерность 9};
 \item тяжёлый модуль переноса: \textbf{размерность 10, инвариантен};
 \item калибровочная орбита связывающего фона: \textbf{размерность 5};
 \item полное замыкание переноса: \textbf{размерность 15};
 \item полное шумовое факторпространство: \textbf{размерность 27};
 \item коммутант стабилизаторной ветви: \textbf{10};
 \item коммутант полной калибровочной ветви: \textbf{16};
 \item следовая изотропия из физической симметрии: \textbf{не следует};
 \item предыдущий 19-мерный Казимир: \textbf{сохранён на стабилизаторе};
 \item следующий гейт: \textbf{родительский гессиан полного
 27-мерного шумового замыкания}.
\end{enumerate}
\end{protocol}

Машинный сертификат сохранён в
\path{s2t_v8_noise_isotropy_symmetry_admission_gate_results.json}.

\begin{thebibliography}{9}
\bibitem{v8-noise-symmetry-holevo}
A.~S.~Holevo,
\emph{Covariant Quantum Dynamical Semigroups: Unbounded Generators},
arXiv:quant-ph/9701037.
\bibitem{v8-noise-symmetry-kibble}
T.~W.~B.~Kibble,
\emph{Symmetry Breaking in Non-Abelian Gauge Theories},
Phys. Rev. 155 (1967) 1554--1561.
\end{thebibliography}